\documentclass[a4paper,11pt]{report}
\usepackage{enumitem}

\usepackage{qj}

\geometry{left=2.2cm,right=2.2cm,top=2.3cm,bottom=2.3cm}
\setlist[enumerate]{label=\arabic*., leftmargin=2em, itemsep=0.8em}

\begin{document}

\begin{center}
    {\Large \bfseries 2026年北京大学强基计划数学试题}
\end{center}

\begin{enumerate}
    \item 已知圆内接四边形 $AB=BC=CD=1$，$\angle B=120^\circ$，求 $AD$ 的长. 

    \item 若复数 $z$ 满足 $|z|=|z-(2+2i)|=3$，求满足题意的所有复数 $z$ 的乘积. 

    \item 求同余方程 $x^2\equiv 1\pmod{30^{2024}}$ 的解的个数. 

    \item 已知 $xy+yz+zx=-4$，求 $|x|+|y|+|z|$ 的最小值. 

    \item 求单位圆上正 $2026$ 边形的一个顶点到另外 $2025$ 个顶点的距离的乘积. 

    \item 若 $x+y=\sqrt{2x+1}+\sqrt{2y+1}$，求 $x+y$ 的最大值. 

    \item 已知 $\dfrac{\sin A+\sqrt{3}\cos A}{\cos A-\sqrt{3}\sin A}=\tan\dfrac{7\pi}{12}$，求 $\sin 2B+2\cos C$ 的取值范围. 

    \item 若定义在实数域上的函数 $f(x)$ 的图像绕原点旋转 $90^\circ$ 后与自身重合，求其不动点的个数. 

    \item 实数 $x,y$ 满足 $x^2+y^2+xy=1$，求 $x-y$ 的最大值. 

    \item 已知椭圆 $\dfrac{x^2}{4}+y^2=1$，点 $P(a,0)$，过 $P$ 作椭圆的两条切线互相垂直，求 $a$ 的值. 

    \item 已知实数 $a,b,c$ 的绝对值均不小于 $1$，$f(x)=x+\dfrac{1}{x}$，若 $f(a)+f(b)+f(c)=0$，求 $a+b+c$ 的最小值. 

    \item 正四面体的棱长为 $1$，求其对棱中点的距离. 

    \item 已知正实数列 $\{a_n\}$ 满足 $a_1=1$，当 $n\geq 2$ 时，$a_n+\dfrac{2}{a_n}=2S_n$（$S_n$ 为前 $n$ 项和），求 $a_{2026}$ 的整数部分. 

    \item 已知复数 $z$ 满足 $z^5=\overline{z}$，求满足该方程的复数 $z$ 的个数. 

    \item 已知集合 $M=\{1,2,\ldots,2026\}$，从 $M$ 中选取四个数组成公差不为 $0$、长度为 $4$ 的等差数列，求这样的等差数列的个数. 

    \item 已知 $p^3+q^4+r^3=-3$，$(p,q,r)$ 为质数数组，求所有满足条件的质数数组 $(p,q,r)$. 

    \item 已知 $p,q$ 为非负整数，二次函数 $f(x)=x^2+px+q$ 有零点；若正数 $a,b$ 满足 $f(a)<f(b)\leq 1.1001f(a)$，求 $b-a$ 的最大值的整数部分. 

    \item 已知方程 $x^4+ax^3+bx^2+ax+1=0$（$a,b\in\mathbb{R}$）的四个根均为正数，求 $b$ 的取值范围. 

    \item 已知椭圆 $\dfrac{x^2}{9}+\dfrac{y^2}{4}=1$，$F$ 为椭圆右焦点，$P$ 为椭圆上动点，求 $|OP|\cdot |OF|$ 的最大值. 

    \item 已知数列 $\{a_n\}$ 满足 $a_0=0$，$a_1=1$；对任意整数 $n\geq 2$，存在 $k\in\{1,2,\ldots,n\}$，使得 $a_n$ 是 $a_{n-1},a_{n-2},\ldots,a_{n-k}$ 这 $k$ 项的算术平均数，求 $a_{2026}-a_{2025}$ 的最小值. 
\end{enumerate}

\end{document}